\worksheettitle
  {M07\(\star\). Чередование в ряду}
  {\modulemark{M07\(\star\)} Усложнение}

\studentnameblock

\begin{notebox}[Два независимых вопроса]
  Номер позиции помогает определить вид объекта, а разность номеров ---
  число промежутков и расстояние. Не смешивайте эти действия.
\end{notebox}

\alternatingtreesfigure

\begin{practicebox}[1. Нечётные и чётные места]
  Первой стоит сосна, затем деревья чередуются.
  \begin{enumerate}[label=\alph*)]
    \item на 17-м месте стоит \answerblank[25mm];
    \item на 28-м месте стоит \answerblank[25mm];
    \item среди первых 20 деревьев берёз \answerblank[15mm];
    \item среди первых 21 дерева сосен \answerblank[15mm].
  \end{enumerate}
\end{practicebox}

\begin{practicebox}[2. Расстояние и вид]
  Сосны и берёзы чередуются, первой стоит сосна. Шаг равен $3$ м.
  Найдите расстояние между 7-м и 18-м деревьями и определите вид каждого.
  \writinglines{3}
\end{practicebox}

\begin{practicebox}[3. Изменённое начало]
  Теперь первой стоит берёза. Какие деревья стоят на 15-м и 24-м местах?
  Сколько сосен среди первых 30 деревьев? \writinglines{3}
\end{practicebox}

\begin{checkpointbox}[Два условия одновременно]
  В ряду чередуются фонарь и дерево, первым стоит фонарь. Расстояние между
  соседними объектами $4$ м. Какой объект расположен на расстоянии $52$ м
  от первого и каким будет номер его позиции? Обоснуйте ответ.
  \writinglines{4}
\end{checkpointbox}
